\documentclass{article} \usepackage[english]{babel} \usepackage[utf8]{inputenc}
\usepackage{amsmath,amssymb} \usepackage{graphicx} \usepackage{hyperref}

\title{Complexity Analysis of the Dutch Merchant Problem} \author{} \date{}

\begin{document}

\maketitle

\begin{abstract} This document presents a computational complexity analysis of
the Dutch Merchant Problem, demonstrating its NP-hardness via a polynomial-time
reduction from the Prize-Collecting Traveling Salesman Problem (PCTSP) and
establishing its NP-Completeness. \end{abstract}

\section{Introduction} \label{sec:introduction}

The Dutch Merchant Problem (DMP) is a combinatorial optimization problem that
models the strategic route and transaction decisions of a maritime merchant
under multiple constraints. This work formally demonstrates that its associated
decision problem is \textbf{NP-Complete}. The proof is structured in two main
parts: (1) establishing its \textbf{NP-hardness} via a polynomial-time reduction
from the Prize-Collecting TSP, and (2) proving its \textbf{membership in the
class NP} by describing a polynomial-time verification algorithm. This
fundamental result implies that no efficient exact algorithm (running in
polynomial time) exists for solving arbitrary-sized instances unless P = NP,
thus justifying the subsequent development of approximation algorithms and
heuristics in the project's later phases.

\section{Proof of NP-Hardness} \label{sec:np-hardness}

\subsection{Formal Definition of the Dutch Merchant Problem (DMP)}
\label{subsec:def_dmp}

An instance of the \textsc{DutchMerchant} problem is defined by the tuple $(G,
r, \mathcal{M}, \mathcal{P}, B, T_{max}, f)$ where:

\begin{itemize} \item $G=(V, E)$ is a complete, undirected graph representing
the port network. $V = \{v_0, v_1, ..., v_n\}$, with $v_0$ being the origin port
(Amsterdam). \item $r \in \mathbb{R}^+$ is the initial capital. \item
$\mathcal{M}$ is the set of $k$ commodity types. \item $\mathcal{P} = \{
(c^{buy}_{v,m}, c^{sell}_{v,m}) \mid v \in V \setminus \{v_0\}, m \in
\mathcal{M} \}$ defines the purchase and sale price for each commodity at each
port. \item $B \in \mathbb{N}$ is the maximum cargo capacity of the ship (in
commodity units). \item $T_{max} \in \mathbb{R}^+$ is the total maximum allowed
duration for the expedition. \item $f: E \rightarrow \mathbb{R}^+$ assigns to
each edge $e=(u,v)$ a travel time $t_{uv}$ and an operational cost $c_{uv}$.
\end{itemize}

A \textbf{feasible solution} is a simple cycle $(v_0, v_{i_1}, ..., v_{i_l},
v_0)$ and a sequence of buy/sell decisions at each visited port, such that: 1)
The capacity $B$ is never exceeded. 2) The capital never becomes negative after
covering operational costs. 3) The total time (travel + operations) $\leq
T_{max}$.

The \textbf{objective} is to maximize the final capital upon returning to $v_0$.

\subsection{Definition of the Reduced Problem: \textsc{Restricted-DMP}}
\label{subsec:def_dmp_rest}

For the reduction, we define a heavily restricted version of
\textsc{DutchMerchant}, called \textsc{Restricted-DMP}. This instance is
obtained by fixing or eliminating the following parameters and decisions from
the original problem:

\begin{enumerate} \item \textbf{Commodities}: We fix $|\mathcal{M}| = 1$. Only
a single generic commodity exists. \item \textbf{Prices and Transactions}: We
fix $c^{buy}_{v} = 0$ and $c^{sell}_{v} = p_v \geq 0$ for every port $v \in V
\setminus \{v_0\}$. This implies that upon visiting $v$, a fixed net profit
$p_v$ is obtained automatically, with no intermediate purchase decision or
capacity consumption. \item \textbf{Capacity}: We fix $B = \infty$. Capacity
ceases to be a constraint. \item \textbf{Financial Constraint}: The intermediate
non-negative capital condition is removed. The capital is simply $r + \sum p_v -
\sum c_{uv}$. \item \textbf{Time}: We fix $T_{max} = \infty$ or to a value
sufficiently large to not restrict any possible tour. The edge costs $c_{uv}$
represent capital *penalties* for travel. \end{enumerate}

Therefore, an instance of \textsc{Restricted-DMP} is defined simply by $(G, v_0,
\{p_v\}_{v \in V}, \{c_{e}\}_{e \in E})$, and its objective is to find a cycle
starting and ending at $v_0$, visiting a subset $S \subseteq V \setminus
\{v_0\}$, that maximizes $\sum_{v \in S} p_v - \sum_{e \in \text{tour}} c_{e}$.

\subsection{Polynomial-Time Reduction from \textsc{Prize-Collecting TSP}}
\label{subsec:reduction_pctsp}

\subsubsection*{Definition of PCTSP} An instance of the
\textsc{PrizeCollectingTSP} (PCTSP) problem is given by $(H, s, \{\pi_u\}_{u \in U}, \{\omega_{ij}\}_{(i,j) \in F})$, where $H=(U, F)$ is a complete graph, $s
\in U$ is the root node, $\pi_u \geq 0$ is the prize for visiting node $u$, and
$\omega_{ij} \geq 0$ is the cost of traversing edge $(i,j)$. The goal is to find
a cycle that starts and ends at $s$, visits a subset of nodes, and maximizes the
sum of collected prizes minus the sum of traversal costs.

\textbf{Theorem 1 (NP-Hardness of PCTSP)}: \textsc{PrizeCollectingTSP} is NP-hard. This follows directly from it being a generalization of the classic \textsc{TravelingSalesmanProblem} (TSP, NP-hard) where the prizes for all nodes except $s$ are greater than the maximum possible sum of travel costs. Furthermore, the NP-hardness of this problem is well-established in the literature, as documented in studies on prize-collecting routing problems \cite{MarchettiSpaccamela2007PrizeCollectingTS}.

\subsubsection*{Construction of the Reduction} Given an arbitrary instance
$I_{P} = (H, s, \{\pi_u\}, \{\omega_{ij}\})$ of PCTSP, we construct an instance
$I_{C} = (G, v_0, \{p_v\}, \{c_{e}\})$ of \textsc{Restricted-DMP} in polynomial
time as follows: \begin{enumerate} \item For each node $u \in U$ in $H$, we
create a corresponding port $v \in V$ in $G$. The root node $s$ is mapped to the
origin port $v_0$. \item The graph $G$ is complete. For each edge $(i,j)$ in $H$
with cost $\omega_{ij}$, we define the travel/operational cost in $G$ as $c_{v_i
v_j} = \omega_{ij}$. \item For each port $v$ corresponding to a node $u \neq s$,
we define the fixed profit $p_v = \pi_u$. \end{enumerate} This construction runs
in $O(|U|^2)$ time.

\subsubsection*{Solution Equivalence and Optimality} There exists a bijective
correspondence between feasible tours of $I_P$ and $I_C$ that preserves the
objective function. Given a tour $T_P = (s, u_{i_1}, ..., u_{i_m}, s)$ for $I_P$
with value $\sum \pi_{u_{i_k}} - \sum \omega_{tour}$, the corresponding tour
$T_C = (v_0, v_{i_1}, ..., v_{i_m}, v_0)$ in $I_C$ yields a final capital of $r
+ \sum p_{v_{i_k}} - \sum c_{tour}$. Since $r$ is constant, maximizing final
capital in $I_C$ is equivalent to maximizing $\sum p_{v_{i_k}} - \sum c_{tour}$,
which is identical to the objective of $I_P$. The inverse argument is analogous.

\textbf{Lemma 1}: There exists a tour in $I_P$ with value $\geq K$ if and only
if there exists a tour in $I_C$ with final capital $\geq r + K$.

\subsubsection*{Reduction Conclusion} Since PCTSP is NP-hard (Theorem 1) and we
have constructed a polynomial-time, optimality-preserving reduction from PCTSP
to \textsc{Restricted-DMP}, we conclude:

\textbf{Theorem 2}: \textsc{Restricted-DMP} is NP-hard.

\subsection{Complexity Preservation Upon Reintroducing Constraints}
\label{subsec:complexity_preservation}

Theorem 2 demonstrates the NP-hardness of a simplified version ($\Pi_{restr}$)
of our original problem ($\Pi_{orig}$). To definitively conclude that
$\Pi_{orig}$ is NP-hard, we must analyze why reintroducing each eliminated
constraint does not transform the problem into an easier one (in P). We
demonstrate that, for each constraint, $\Pi_{restr}$ is a \textbf{special case}
of $\Pi_{orig}$ under a specific configuration; therefore, $\Pi_{orig}$ is at
least as hard.

Let $\Pi_{orig}(X, Y, Z, ...)$ be the general problem with constraints $X, Y, Z,
...$, and $\Pi_{restr}$ the problem without them. If $\Pi_{restr}$ is NP-hard,
and there exists an assignment of values for $X, Y, Z, ...$ such that
$\Pi_{orig}(val_X, val_Y, val_Z, ...) \equiv \Pi_{restr}$, then $\Pi_{orig}$ is
NP-hard.

\begin{enumerate} \item \textbf{Multiple Commodities ($|\mathcal{M}| \geq 1$)}.
In $\Pi_{restr}$, $|\mathcal{M}|=1$. Consider an instance of $\Pi_{orig}$ with
$|\mathcal{M}|=k>1$ where: i) the purchase prices for commodities $2, ..., k$
are higher than any possible profit, making them non-profitable; ii) the sale
price of commodity 1 is identical to $p_v$; and iii) capacity is sufficient to
carry only commodity 1. Any optimal solution will ignore commodities $2, ..., k$
and behave as in $\Pi_{restr}$. Thus, $\Pi_{restr}$ is a special case of
$\Pi_{orig}$ with multiple commodities, which does not reduce complexity.

\item \textbf{Finite Cargo Capacity ($B \in \mathbb{N}$)}. In $\Pi_{restr}$,
$B = \infty$. Setting $B$ to a sufficiently large value (e.g., $B \geq \sum_{v
\in V} 1$) in $\Pi_{orig}$ replicates the unlimited capacity condition, as the
limit will never be reached with a single unit-profit commodity. Reducing $B$
restricts the solution space, adding a dimension of difficulty equivalent to the
Knapsack problem, which is NP-hard. Adding an NP-hard constraint cannot simplify
a problem already proven NP-hard.

\item \textbf{Intermediate Financial Constraint (Non-Negative Capital)}. In
$\Pi_{restr}$ this constraint is omitted. In $\Pi_{orig}$, if the initial
capital $r$ is greater than the maximum possible sum of travel costs, the
constraint never activates, replicating $\Pi_{restr}$. Reducing $r$ introduces a
dynamic financial state that must be verified at each port, increasing the
complexity of the state space, not reducing it.

\item \textbf{Strict Time Limit ($T_{max} \in \mathbb{R}^+$)}. In
$\Pi_{restr}$, $T_{max} = \infty$. Setting $T_{max}$ larger than the duration of
the longest possible tour in $\Pi_{orig}$ renders the constraint inactive,
replicating $\Pi_{restr}$. Imposing a strict time limit transforms the problem
into a variant of the "Orienteering Problem" (also NP-hard), adding an
additional layer of difficulty.

\item \textbf{Complex Buy/Sell Decisions}. In $\Pi_{restr}$, profit is
automatic ($c^{buy}_v = 0, c^{sell}_v = p_v$). In $\Pi_{orig}$, fixing prices
identically and defining that profit is obtained upon port arrival (without
capacity consumption) simulates $\Pi_{restr}$. Allowing variable prices and
strategic decisions introduces an arbitrage and inventory management problem,
which is a strict and more difficult generalization. \end{enumerate}

\textbf{Analysis Conclusion}: For each constraint $R_i$, there exists a
configuration of $\Pi_{orig}$ that makes it equivalent to $\Pi_{restr}$. Any
adjustment that genuinely activates $R_i$ (e.g., reducing $B$, reducing
$T_{max}$) \textbf{restricts} the feasible solution space and/or adds an NP-hard
optimization dimension. Therefore, $\Pi_{orig}$ is no simpler than $\Pi_{restr}$;
it is a generalization that preserves and potentially increases complexity.

\textbf{Theorem 3 (NP-Hardness of the General Problem)}: The
\textsc{DutchMerchant} problem ($\Pi_{orig}$) is NP-hard. \textbf{Proof}: Direct
from Theorem 2 and the above analysis, which establishes that $\Pi_{restr}$ is a
special case of $\Pi_{orig}$. \hfill $\square$

\subsection{Proof of Membership in NP} \label{subsec:membership_np}

To establish that \textsc{DutchMerchant} is NP-Complete, in addition to being
NP-hard (Theorem 3), we must prove it belongs to the class NP. That is, given a
proposed solution to the problem, we can verify its feasibility and compute its
objective value (the final capital) in \textbf{polynomial time} relative to the
input size.

\subsubsection*{Certificate and Verifier} A certificate $C$ for an instance of
$\Pi_{orig}$ must include: \begin{enumerate} \item A sequence of ports (a tour)
$T = (v_0, v_{i_1}, v_{i_2}, ..., v_{i_m}, v_0)$. \item For each visited port
$v_{i_j}$, a list of buy and sell operations per commodity type. \end{enumerate}

The size of this certificate is polynomially bounded: the tour has at most
$|V|$ nodes, and operations at each port are for $|\mathcal{M}|$ commodities.

A verifier algorithm $V(I, C)$ must check the following points in polynomial
time: \begin{enumerate} \item \textbf{Tour Validity}: Verify $T$ is a cycle
starting and ending at $v_0$ and that no internal port is repeated. This takes
$O(|V|)$. \item \textbf{Capacity Compliance}: For each leg between ports,
calculate the total onboard load by summing purchases and subtracting sales.
Verify it never exceeds $B$. This requires traversing the operation sequence
once: $O(m \cdot |\mathcal{M}|)$. \item \textbf{Financial Constraint Compliance}:
Simulate the cash flow. Starting from $r$, for each port: subtract travel costs,
add sale profits, subtract purchase costs, and verify non-negative capital.
Also $O(m \cdot |\mathcal{M}|)$. \item \textbf{Time Limit Compliance}: Sum
travel times $t_{uv}$ for all edges in $T$ and port operation times. Verify the
total does not exceed $T_{max}$. $O(m)$. \item \textbf{Objective Calculation}:
The certificate's value is the final capital after the last operation, obtained
as the output of step 3. \end{enumerate}

Each operation is an addition, subtraction, or comparison executed a number of
times proportional to the certificate size. Therefore, total verification time
is polynomial in the input size $I$ (which includes $|V|$, $|E|$,
$|\mathcal{M}|$).

\textbf{Lemma 3 (Membership in NP)}: There exists a verifier $V(I, C)$ that
decides in polynomial time whether a certificate $C$ is a feasible solution for
$\Pi_{orig}$ and what its value is. Therefore, $\Pi_{orig} \in \text{NP}$.

\subsubsection*{NP-Completeness Conclusion} Combining the NP-hardness result
(Theorem 3) with NP membership (Lemma 3), we establish the final classification.

\textbf{Theorem 4 (NP-Completeness)}: The decision problem associated with
\textsc{DutchMerchant} is NP-Complete. \textbf{Proof}: The decision problem asks:
"Does there exist a route and transaction plan with final capital $\geq K$?".
By Lemma 3, this problem is in NP. By Theorem 3, it is NP-hard. Consequently, it
is NP-Complete. \hfill $\square$

This classification formally justifies the use of non-exact algorithmic
techniques (heuristics, approximations) in the project's subsequent phases.

\section{Conclusions} \label{sec:conclusions}

This work presents a formal analysis of the computational complexity of the
Dutch Merchant Problem. The main results are: \begin{itemize} \item A restricted
version of the problem (\textsc{Restricted-DMP}) was defined and proven NP-hard
via a polynomial-time reduction from the Prize-Collecting TSP (Theorem 2). \item
It was rigorously argued, constraint by constraint, that the NP-hard complexity
is preserved when reintroducing all features of the general problem (Theorem 3).
\item It was demonstrated that the associated decision problem belongs to the
class NP, by presenting a polynomial-time verifiable certificate (Lemma 3 and
Theorem 4). \end{itemize} Therefore, we conclude that the \textbf{Dutch Merchant
decision problem is NP-Complete}. This finding provides the theoretical
justification for the project's practical approach, which in the following
phases will focus on the design, implementation, and empirical evaluation of
approximation algorithms and heuristics to tackle realistic problem instances.

\bibliographystyle{plain} \bibliography{referencias}

\end{document}